\bookchapter{第九章：モバイルオフィスと学習}

\chapterintro{一、ドキュメント処理：掌中での編集と創作}

\booksection{スマホ画面でドキュメントを編集するのは、パソコンとは確かに感覚が違います。最大の問題は画面が小さく、指が太く、レイアウトがずれやすいこと。でもいくつかのコツで体験はかなり良くなります：}

WPS Office

\begin{itemize}
  \item \textbf{：ローカライズが最も充実しており、Microsoft形式（docx/xlsx/pptx）と完全互換、内蔵のPDFツールセットが豊富。私が特に気に入っているのは「音声メモ」機能です。契約書を見る時にタイピングしにくい場合、スマホに向かって意見を数言話すだけで、自動的にテキスト注釈になってドキュメントに貼り付けられます。またWPSには大量の無料テンプレートが内蔵されており、履歴書や報告PPTを作る時に時間を大幅に節約できます。}Microsoft 365
  \item \textbf{：会社が正規のOfficeを使っているなら、これがベストチョイス。デスクトップ版とのフォーマット互換性はほぼ完璧で、WPSで複雑なレイアウトを開くとわずかにずれるようなことがありません。OneDriveと深く統合されており、パソコンで編集途中のドキュメントも、スマホで開けば最新版になっています。}石墨ドキュメント/テンセントドキュメント
  \item \textbf{：この2つは「オンライン協業」の代表。複数人で同時に一つのドキュメントを編集するのに適しており、各自の変更がリアルタイムで表示されます。私は石墨ドキュメントを使って友人と旅行ガイドを一緒に作りました。二人で同時に表に情報を入力する効率は、一人でやるよりはるかに高いです。}実用的なワークフロー
\end{itemize}

\textbf{：私の現在の習慣は、重要なファイルはOneDriveに保存、日常的な編集はWPS、チーム協業は石墨ドキュメントを使う。ファイルをクラウドに置くようになってから、「USBメモリを忘れた」という恥ずかしい思いはもうありません。}二、ノートアプリ：手軽なメモから知識体系へ

\booksection{ノートを取るという事について、多くの人は最初から複雑なシステムを作ろうとして、三日も続きません。私の提案は、まず軽く始めて重くする。「記録する」習慣を身につけてから、「管理する」問題を考える。}

手軽なメモ——システム標準で十分

\begin{itemize}
  \item \textbf{：iOSメモは起動が極めて速く、書類スキャン、手書き、録音にも対応しています。Android各社のメモ帳も劣らず、ロック画面状態から素早く起動できます。これらのツールは一時的なアイデア、買い物リスト、会議の要点を記録するのに適しています。重要なのは起動が速く、思いついたらすぐ記録でき、Appの読み込みを待つ必要がないこと。}後で読む——情報収集
  \item \textbf{：WeChatの「フローティングウィンドウ」機能は実は隠れた神器です。良い記事を見てじっくり読む時間がない時、右上の「フローティングウィンドウ」をタップすると、小さな円いアイコンになって画面の端に貼り付けられ、いつでもタップして読み続けられます。「ファイル転送アシスタント」に転送するよりずっと便利。また、CuboxやPocketのような「後で読む」ツールは、どこで見た素晴らしいコンテンツでもクロスプラットフォームで収集できます。}体系構築——Notionと飛書
  \item \textbf{：散在するノートが探しにくくなってきたら、Notionに移行しましょう。これは本当に柔軟性が高い——「読書ノートデータベース」を作って、書名、著者、評価で自動的に整理できますし、「プロジェクトカンバン」を作って作業タスクの進捗を可視化することもできます。飛書の「多次元表」も同様で、かつ国内からのアクセス速度が速いです。私はNotionですべての執筆素材、学習ノート、プロジェクトの進捗を管理しており、もう3年以上使っています。}実用的なノート戦略
\end{itemize}

\textbf{：私自身の習慣は「三種の神器」——一時的なアイデアはシステムメモ帳（秒で起動）、後で読む記事はWeChatフローティングウィンドウ（手軽）、長期蓄積が必要な知識はNotion（体系化）。最初からすべてを一つのツールに詰め込もうとせず、そうすると逆に挫折しやすくなります。}三、スキャンとOCR：紙の世界をスマホに取り込む

\booksection{スマホスキャンは、うまく使えば印刷屋に走る手間をかなり省けます。前回ビザ申請で大量の書類の電子版が必要になった時、すべてスマホで対応しました。}

iPhoneユーザー

\begin{itemize}
  \item \textbf{：iOSメモ標準のスキャン機能がすでに十分優れています。メモを開く カメラアイコンをタップ 「書類をスキャン」を選択、書類に向けると自動的に枠を認識して透視を補正します。撮影後に色調整ができます——「白黒」モードを選ぶとスキャナーに最も近い効果になり、「グレースケール」は画像入りの資料に適しています。}Androidユーザー
  \item \textbf{Appleユーザー}OCR文字認識
  \item \textbf{：スキャンした書類の文字を抽出して編集するには、OCRが必要です。WeChat標準の「文字抽出」機能が最も簡単——画像を長押しすれば使え、認識精度もなかなかです。大量処理が必要な場合は、白描Appや全能掃描王を使いましょう。中国語認識では、BaiduのOCRエンジンの効果が比較的良いです。WPSの画像からテキストへの変換機能も精度が高いです。}スキャンの小技
  \item \textbf{：明るい場所で撮影し、できればデスクライトを横から当てて、書類に手の影が落ちないようにします。書類の角が折れている場合は、手で平らにしてから撮りましょう。複数ページの書類は連続スキャン後に一つのPDFに統合でき、ページごとに別々に保存する必要はありません。}四、クラウドストレージの選び方：ベストはない、あなたに最適なものだけ
\end{itemize}

\booksection{クラウドストレージ選びについて私は長く考えましたが、結局「メインのエコシステム」が何かによるという結論に達しました。}

Appleエコシステム iCloud

\begin{itemize}
  \item \textbf{：その存在を感じさせないほどシームレスに同期されます。スマホで撮った写真が、数分後にはパソコンに自動的に入っています。またAppleのプライバシー保護は最も厳格です。ただし欠点も明らかで、AndroidやWindowsユーザーへのクロスプラットフォーム共有が面倒、無料枠は5GBのみ、拡張には費用がかかります。}Windowsメイン OneDrive
  \item \textbf{：WindowsパソコンとOffice 365を使っているなら、OneDriveはまさに天作の組み合わせ。Wordでドキュメントを保存すると自動的にクラウドに保存されます。1TBの容量がOfficeサブスクに付属しておりお得です。スマホAppの自動写真バックアップ機能も使いやすいです。}国内共有 Baiduネットディスク
  \item \textbf{：正直、Baiduネットディスクのダウンロード速度制限はずっと批判されてきましたが、その共有エコシステムは確かに他の追随を許しません。中国語インターネット上の大量の学習資料や映像素材はBaiduネットディスクで共有されています。たまに他人が共有したファイルをダウンロードするだけなら、会員にならなくてもなんとかなります。ヘビーな資源収集家なら、会員になった方がいいでしょう。節約した時間の方が価値があります。}Alibabaクラウドドライブ/Quarkネットディスク
  \item \textbf{：この2年で台頭してきた新鋭。Alibabaクラウドドライブはアップロード·ダウンロード無制限（現時点）、QuarkネットディスクはQuarkブラウザと深く連携しており、資料検索時に直接ネットディスクに保存できて便利です。}私の組み合わせ
\end{itemize}

\textbf{：iCloudで個人写真と日常ファイル（プライバシー優先）、Baiduネットディスクで資源保存と共有受信（エコシステム優先）、OneDriveで仕事ドキュメント（Office統合優先）。3つのクラウドストレージがそれぞれの役割を担い、干渉しません。}五、実際のモバイルオフィスシーン

\booksection{私自身の実体験をお話ししましょう。ある時出張先で、クライアントから急に契約条項の修正を求められました。私のワークフローはこうでした：}

ホテルでクライアントからWeChatで修正意見を受け取る（テキストメッセージ＋手書き注釈の写真1枚）。

\begin{enumerate}
  \item OneDriveを開いて契約書原本を探し、WPSで開いて編集、クライアントのWeChatメッセージを参照しながら条項を一つずつ修正。
  \item 修正完了後、WPSの「PDFとして出力」機能で書き出し、クライアントのWeChatグループに送って確認。
  \item クライアントから第二ラウンドの意見が出され、WPSのペンツールで直接PDF上に修正箇所を丸で囲み、音声注釈で修正の意図を説明。
  \item 確認OK後、WPSの「電子署名」機能でオンライン署名し、相手にも同様に署名してもらう。
  \item 最後に署名済みの契約書PDFをBaiduネットディスクとiCloudに1部ずつ保存し、同時にNotionの作業ログに記録。
  \item 全過程は開始から終了まで約40分で、パソコンは一切開きませんでした。正直、3年前はスマホでこれらができるとは信じませんでしたが、今では確実にできます。重要なのは、これらのツールに事前に慣れておくこと。必要な時に慌ててダウンロード·登録するのでは遅すぎます。
\end{enumerate}

やってみよう：モバイルオフィスのワークフローを構築する

\begin{tcbpractice}{メインのドキュメントApp（WPSまたはMicrosoft 365）を選び、ログインしてクラウドストレージを連携。}
\begin{enumerate}
  \item メモまたは付箋を開き、スキャン機能で書類を撮影し、PDFとして書き出して自分に送信。
  \item クラウドストレージAppを一つ選び、スマホ写真の自動バックアップに設定（「Wi-Fi時のみバックアップ」で通信量節約）。
  \item Notionまたは飛書で新規ページを作成し、簡単な「作業ログ」テンプレートを構築してみる。
  \item 第十章：モバイル決済と生活サービス
\end{enumerate}

\end{tcbpractice}
